\section*{Chapitre \ref{ch:m1:what-a-trading-firm-does} --- Ce que fait une société de trading}

\begin{solution}{exo:m1:what-a-trading-firm-does:1}
\omterm{def:m1:what-a-trading-firm-does:spread}{Fourchette} $25.02 - 24.98 = \$0.04$~; \omterm{def:m1:what-a-trading-firm-does:spread}{prix milieu} $\$25.00$~; $0.04/25.00 = \qty{16}{\bp}$.
Un ordre au marché paie la \omterm{def:m1:what-a-trading-firm-does:spread}{demi-fourchette}, $\$0.02$ par action~: $300 \times 0.02 =
\$6.00$.
\end{solution}

\begin{solution}{exo:m1:what-a-trading-firm-does:2}
$2\,000\,000 \times \$0.001 = \$2\,000$ par jour, et $\$504\,000$ par an~: un
dixième de cent n'est pas une activité rentable à ce volume, c'est tout l'intérêt de
la \cref{prop:m1:what-a-trading-firm-does:breakeven}.
\end{solution}

\begin{solution}{exo:m1:what-a-trading-firm-does:3}
(a) \omterm{def:m1:what-a-trading-firm-does:principal-agent}{Agent}~; commission. (b) \omterm{def:m1:what-a-trading-firm-does:principal-agent}{Principal}~; \omterm{def:m1:what-a-trading-firm-does:spread}{fourchette} (la décote), avec risque de position.
(c) \omterm{def:m1:what-a-trading-firm-does:principal-agent}{Agent} au sens économique --- le client possède le portefeuille~; commission.
(d) \omterm{def:m1:what-a-trading-firm-does:principal-agent}{Principal}~; \omterm{def:m1:what-a-trading-firm-does:spread}{fourchette}. (e) \omterm{def:m1:what-a-trading-firm-does:principal-agent}{Principal} pour le compte des investisseurs du fonds~; \omterm{def:m1:what-a-trading-firm-does:alpha}{alpha} --- et,
si le modèle se contente de charger sur le risque de marché, prime de risque.
\end{solution}

\begin{solution}{exo:m1:what-a-trading-firm-does:4}
Profit brut $0.12 \times 500 = 60$ millions de dollars. Frais de gestion 10~millions de dollars~;
commission de performance $0.20 \times (60 - 10) = 10$ millions de dollars. Les investisseurs gardent
40~millions de dollars~: un rendement net de 8~\%. Le gestionnaire perçoit $20/60$, soit un tiers
du profit brut.
\end{solution}

\begin{solution}{exo:m1:what-a-trading-firm-does:5}
$c = 1.00 - 0.14 \times 5 + 0.20 - 0.05 = 0.45$ cent. $V^\star = 90\,000 /
0.0045 = 20$ millions d'actions par jour.
\end{solution}

\begin{solution}{exo:m1:what-a-trading-firm-does:6}
(a) $467.8 \times 10^6 / (1\,545 \times 10^9) = \qty{3.03}{\bp}$. C'est une
borne supérieure, car le revenu net de négociation inclut aussi les opérations de la société sur
d'autres classes d'actifs, alors que le dénominateur ne retient que les produits
cotés. (b) $467.8/609 = \text{\euro}0.77$ million par salarié. (c)
$50 \times 3.03 \times 10^{-4} = \text{\euro}0.015$~: un centime et demi.
\end{solution}

\begin{solution}{exo:m1:what-a-trading-firm-does:7}
(a) \omterm{def:m1:what-a-trading-firm-does:spread}{Fourchette} gagnée 20\,000~\$~; résultat de position $-\$9\,173$~; net 14\,427~\$ (après
4\,000~\$ de rabais et 400~\$ de compensation). (b) La dispersion provient de
l'inventaire~: la position du \omterm{def:m1:what-a-trading-firm-does:market-maker}{teneur de marché} est une marche aléatoire de $\pm$ une taille
d'ordre par transaction, et cette position est exposée à chaque mouvement ultérieur du
\omterm{def:m1:what-a-trading-firm-does:spread}{prix milieu}, sauts informés compris. Quadrupler la taille des ordres tout en divisant le
nombre d'ordres par quatre conserve le volume mais rend la marche de l'inventaire
$4\sqrt{1/4} = 2$ fois plus grande à tout moment de la journée~: l'écart-type
du résultat quotidien double à peu près, à moyenne inchangée.
\end{solution}

\begin{solution}{exo:m1:what-a-trading-firm-does:8}
D'abord, les ordres au marché \emph{paient} la \omterm{def:m1:what-a-trading-firm-does:spread}{demi-fourchette}~; seuls les ordres passifs la perçoivent,
et leur exécution n'est pas garantie~: la stratégie telle que décrite perd 6~points
de base par aller-retour. Ensuite, même une stratégie passive ne conserve pas la
\omterm{def:m1:what-a-trading-firm-does:spread}{fourchette}~: la \omterm{def:m1:what-a-trading-firm-does:adverse-selection}{sélection adverse} en restitue une large part
(la \cref{prop:m1:what-a-trading-firm-does:capture}), et les exécutions sont plus probables
précisément lorsqu'elles sont déficitaires. Avec dix allers-retours par jour au marché, le
chiffre probable est d'environ $-60$ points de base par jour avant frais.
\end{solution}

\begin{solution}{pb:m1:what-a-trading-firm-does:1}
\textbf{1.} 1 cent~; $0.01/40 = \qty{2.5}{\bp}$.
\textbf{2.} $0.12 \times 5 = 0.60$ cent.
\textbf{3.} $c = 1.00 - 0.60 + 0.20 - 0.02 = 0.58$ cent.
\textbf{4.} $6 \times 10^6 \times 0.0058 = \$34\,800$.
\textbf{5.} $\$8\,769\,600$.
\textbf{6.} $4.9 + 2.1 + 0.56 = 7.56$ millions de dollars~; 30\,000~\$ par jour.
\textbf{7.} 4\,800~\$ par jour~; 1\,209\,600~\$ par an.
\textbf{8.} $30\,000/34\,800 = 86.2\%$.
\textbf{9.} $V^\star = 30\,000/0.0058 = 5\,172\,414$ actions~: la société se situe 16~\%
au-dessus du seuil de rentabilité.
\textbf{10.} $c = 1.00 - 0.80 + 0.18 = 0.38$ cent~; $V^\star = 7\,894\,737$.
\textbf{11.} $6 \times 10^6 \times 0.0038 - 30\,000 = -\$7\,200$ par jour.
\textbf{12.} Une \omterm{def:m1:what-a-trading-firm-does:spread}{demi-fourchette} de 1.2 cent, soit une \omterm{def:m1:what-a-trading-firm-does:spread}{fourchette} de 2.4 cents, qui ne
peut pas être cotée. Quad Street peut coter 3 cents et perdre l'essentiel de ses
exécutions au profit de concurrents restés à 2, ou rester à 2 et ne coter que lorsque ses propres signaux
indiquent que le flux est favorable~: choisir \emph{quand} être présent sur le marché est le seul
contrôle continu dont elle dispose.
\textbf{13.} $c = 0.53$ cent~; $V^\star = 5\,660\,377$.
\textbf{14.} $\Phi(-4\,800/22\,000) = \Phi(-0.218) = 0.41$.
\textbf{15.} $0.218 \times \sqrt{252} = 3.46$.
\textbf{16.} $\Phi(-3.46) = 0.027\%$~: environ une année sur 3\,700.
\textbf{17.} $1\,209\,600/5\,000\,000 = 24.2\%$.
\textbf{18.} $(69\,600 - 30\,000)/4\,800 = 8.25$.
\textbf{19.} Les marchés volatils apportent du volume et des \omterm{def:m1:what-a-trading-firm-does:spread}{fourchettes} plus larges~; avec
des coûts fixes, presque tout le revenu supplémentaire est un profit (question 18). Ils
entraînent aussi des variations plus fortes contre l'inventaire et davantage de flux informés, de sorte que l'écart-type
du résultat quotidien et le terme de \omterm{def:m1:what-a-trading-firm-does:adverse-selection}{sélection adverse} augmentent avec eux.
\textbf{20.} Quad Street survit au-delà d'environ 5.2 millions d'actions par jour, et
les trois chiffres à suivre sont son volume, sa \omterm{prop:m1:what-a-trading-firm-does:capture}{captation nette} par action
(surtout la part liée à la \omterm{def:m1:what-a-trading-firm-does:adverse-selection}{sélection adverse}) et son coût par jour.
\end{solution}

\begin{solution}{iq:m1:what-a-trading-firm-does:1}
Elle gagne la \omterm{def:m1:what-a-trading-firm-does:spread}{demi-fourchette} sur chaque exécution, plus tout rabais de bourse, en vendant
de l'immédiateté à ceux qui veulent négocier tout de suite. Elle perd de deux façons~: la \omterm{def:m1:what-a-trading-firm-does:adverse-selection}{sélection adverse}
--- les exécutions obtenues sont de manière disproportionnée celles qui précèdent immédiatement le
mouvement défavorable du prix --- et le risque d'inventaire sur la position qu'elle accumule
entre-temps. Elle supporte des coûts fixes de technologie et de personnel, de sorte que le profit est une
faible \omterm{prop:m1:what-a-trading-firm-does:capture}{captation nette} multipliée par un grand volume, diminuée de ces coûts.
\iqlookfor{les deux pertes nommées sans aide, et le sentiment que la \omterm{def:m1:what-a-trading-firm-does:spread}{fourchette}
est un chiffre brut, et non le profit.}
\end{solution}

\begin{solution}{iq:m1:what-a-trading-firm-does:2}
Une \omterm{def:m1:what-a-trading-firm-does:prop}{société de trading pour compte propre} négocie l'argent de ses propriétaires et tire son profit de négociation~; un
\omterm{def:m1:what-a-trading-firm-does:hedge-fund}{hedge fund} négocie l'argent des investisseurs et perçoit des frais sur les encours et sur les profits.
Conséquences~: le fonds recherche la \emph{capacité} (les frais croissent avec les encours) et
des rendements mensuels réguliers que les investisseurs accepteront~; la \omterm{def:m1:what-a-trading-firm-does:prop}{société de trading pour
compte propre} vise un rendement sur une faible base de capital et peut mener des stratégies à faible
capacité, des résultats plus volatils, et aucune obligation de les expliquer à quiconque.
\iqlookfor{incitations découlant de qui possède le capital, et non d'une liste de
noms de sociétés.}
\end{solution}

\begin{solution}{iq:m1:what-a-trading-firm-does:3}
On achète 100 à 49.99 et le \omterm{def:m1:what-a-trading-firm-does:spread}{prix milieu} vaut désormais 49.97~: $-\$2$. Décomposition~: $+\$1$ de
\omterm{def:m1:what-a-trading-firm-does:spread}{fourchette} gagnée (on a acheté 1 cent sous le \omterm{def:m1:what-a-trading-firm-does:spread}{prix milieu} de 50.00 alors en vigueur) et $-\$3$ dû à
la baisse de 3 cents du \omterm{def:m1:what-a-trading-firm-does:spread}{prix milieu} par la suite. La seconde partie, moyennée sur de nombreuses
exécutions, constitue la \omterm{def:m1:what-a-trading-firm-does:adverse-selection}{sélection adverse} subie.
\iqlookfor{valorisation au \omterm{def:m1:what-a-trading-firm-does:spread}{prix milieu}, et décomposition proposée spontanément.}
\end{solution}

\begin{solution}{iq:m1:what-a-trading-firm-does:4}
Un point de base sur 2~milliards de dollars représente 200\,000~\$ par jour~; multiplié par environ 250 jours,
50~millions de dollars par an. Le montant est plausible~: les chiffres publiés par les \omterm{def:m1:what-a-trading-firm-does:market-maker}{teneurs de marché}
cotés impliquent au plus quelques points de base de valeur échangée, et beaucoup moins dans
les produits les plus liquides.
\iqlookfor{calcul mental rapide ($2 \times 10^9 \times 10^{-4}$) et contrôle
de vraisemblance par rapport à une donnée connue.}
\end{solution}

\begin{solution}{iq:m1:what-a-trading-firm-does:5}
En \omterm{def:m1:what-a-trading-firm-does:principal-agent}{agent}~: exécuter l'ordre sur plusieurs jours avec un algorithme d'exécution, moyennant
une commission~; le client conserve le risque de prix et le risque de fuite
d'information. En \omterm{def:m1:what-a-trading-firm-does:principal-agent}{principal}~: se porter acheteur du bloc à une décote par rapport au dernier cours~; le
desk prend le risque et est rémunéré par la décote, qui doit couvrir le
coût attendu de la liquidation de dix jours de volume, plus la volatilité sur
cette période. Des formules hybrides existent~: un prix garanti calé sur un indice de référence pour une partie, et
l'agence pour le reste.
\iqlookfor{le transfert de risque explicitement énoncé, et une décote qui croît
avec la taille et la volatilité.}
\end{solution}

\begin{solution}{iq:m1:what-a-trading-firm-does:6}
(i) Les \omterm{def:m1:what-a-trading-firm-does:spread}{fourchettes} se sont élargies \emph{parce que} le flux d'ordres est devenu plus informé ou plus
volatil~: la \omterm{def:m1:what-a-trading-firm-does:adverse-selection}{sélection adverse} par action a augmenté de plus que la \omterm{def:m1:what-a-trading-firm-does:spread}{demi-fourchette}.
(ii) La \omterm{def:m1:what-a-trading-firm-does:spread}{fourchette} était plus large que celle des concurrents~: les ordres n'étaient exécutés
que lorsque tous les autres s'étaient retirés --- les pires exécutions --- et que le volume était
entretenu justement par ces exécutions-là. Mesure~: le markout des exécutions (le \omterm{def:m1:what-a-trading-firm-does:spread}{prix milieu} quelques
secondes et minutes après chaque exécution, signé selon le côté de l'ordre), ventilé selon que la
société était seule au meilleur cours.
\iqlookfor{l'idée que les exécutions sont sélectionnées, et une mesure concrète
plutôt qu'une histoire.}
\end{solution}
